%beamer

% Comment/uncomment this line to toggle handout mode
% \newcommand{\handout}{}

% % \bigtimes abgeschrieben von http://tex.stackexchange.com/questions/14386/importing-a-single-symbol-from-a-different-font
% \DeclareFontFamily{U}{mathx}{\hyphenchar\font45}
% \DeclareFontShape{U}{mathx}{m}{n}{
%       <5> <6> <7> <8> <9> <10> gen * mathx
%       <10.95> mathx10 <12> <14.4> <17.28> <20.74> <24.88> mathx12
%       }{}
% \DeclareSymbolFont{mathx}{U}{mathx}{m}{n}
% \DeclareMathSymbol{\bigtimes}{\mathop}{mathx}{161}

\RequirePackage{xcolor}

\def\9{\square}
%\def\9{\blank}

% f"ur Aussagenlogik
\colorlet{alcolor}{blue}
\RequirePackage{tikz}
\usetikzlibrary{arrows.meta}
\newcommand{\alimpl}{\mathrel{\tikz[x={(0.1ex,0ex)},y={(0ex,0.1ex)},>={Classical TikZ Rightarrow[]}]{\draw[alcolor,->,line width=0.7pt,line cap=round] (0,0) -- (15,0);\path (0,-6);}}}
\newcommand{\aleqv}{\mathrel{\tikz[x={(0.1ex,0ex)},y={(0ex,0.1ex)},>={Classical TikZ Rightarrow[]}]{\draw[alcolor,<->,line width=0.7pt,line cap=round] (0,0) -- (18,0);\path (0,-6);}}}
\newcommand{\aland}{\mathbin{\raisebox{-0.6pt}{\rotatebox{90}{\texttt{\color{alcolor}\char62}}}}}
\newcommand{\alor}{\mathbin{\raisebox{-0.8pt}{\rotatebox{90}{\texttt{\color{alcolor}\char60}}}}}
%\newcommand{\ali}[1]{_{\mathtt{\color{alcolor}#1}}}
\newcommand{\alv}[1]{\mathtt{\color{alcolor}#1}}
\newcommand{\alnot}{\mathop{\tikz[x={(0.1ex,0ex)},y={(0ex,0.1ex)}]{\draw[alcolor,line width=0.7pt,line cap=round,line join=round] (0,0) -- (10,0) -- (10,-4);\path (0,-8) ;}}}
\newcommand{\alP}{\alv{P}} %ali{#1}}
%\newcommand{\alka}{\negthinspace\hbox{\texttt{\color{alcolor}(}}}
\newcommand{\alka}{\negthinspace\text{\texttt{\color{alcolor}(}}}
%\newcommand{\alkz}{\texttt{\color{alcolor})}}\negthinspace}
\newcommand{\alkz}{\text{\texttt{\color{alcolor})}}\negthinspace}
\newcommand{\AAL}{A_{AL}}
\newcommand{\LAL}{\hbox{\textit{For}}_{AL}}
\newcommand{\AxAL}{\hbox{\textit{Ax}}_{AL}}
\newcommand{\AxEq}{\hbox{\textit{Ax}}_{Eq}}
\newcommand{\AxPL}{\hbox{\textit{Ax}}_{PL}}
\newcommand{\AALV}{\hbox{\textit{Var}}_{AL}}
\newcommand{\MP}{\hbox{\textit{MP}}}
\newcommand{\GEN}{\hbox{\textit{GEN}}}
\newcommand{\W}{\ensuremath{\hbox{\textbf{w}}}\xspace}
\newcommand{\F}{\ensuremath{\hbox{\textbf{f}}}\xspace}
\newcommand{\WF}{\ensuremath{\{\W,\F\}}\xspace}
\newcommand{\val}{\hbox{\textit{val}}}
\newcommand{\valDIb}{\val_{D,I,\beta}}

\newcommand*{\from}{\colon}

% die nachfolgenden Sachen angepasst an cmtt
\newlength{\ttquantwd}
\setlength{\ttquantwd}{1ex}
\newlength{\ttquantht}
\setlength{\ttquantht}{6.75pt}
\def\plall{%
  \tikz[line width=0.67pt,line cap=round,line join=round,baseline=(B),alcolor] {
    \draw (-0.5\ttquantwd,\ttquantht) -- node[coordinate,pos=0.4] (lll){} (-0.25pt,-0.0pt) -- (0.25pt,-0.0pt) -- node[coordinate,pos=0.6] (rrr){} (0.5\ttquantwd,\ttquantht);
    \draw (lll) -- (rrr);
    \coordinate (B) at (0,-0.35pt);
  }%
}
\def\plexist{%
  \tikz[line width=0.67pt,line cap=round,line join=round,baseline=(B),alcolor] {
    \draw (-0.9\ttquantwd,\ttquantht) -- (0,\ttquantht) -- node[coordinate,pos=0.5] (mmm){} (0,0) --  (-0.9\ttquantwd,0);
    \draw (mmm) -- ++(-0.75\ttquantwd,0);
    \coordinate (B) at (0,-0.35pt);
  }\ensuremath{\,}%
}
\let\plexists=\plexist
\newcommand{\NT}[1]{\ensuremath{\langle\mathrm{#1} \rangle}}

\newcommand{\CPL}{\text{\itshape Const}_{PL}}
\newcommand{\FPL}{\text{\itshape Fun}_{PL}}
\newcommand{\RPL}{\text{\itshape Rel}_{PL}}
\newcommand{\VPL}{\text{\itshape Var}_{PL}}
\newcommand{\ATer}{A_{\text{\itshape Ter}}}
\newcommand{\ARel}{A_{\text{\itshape Rel}}}
\newcommand{\AFor}{A_{\text{\itshape For}}}
\newcommand{\LTer}{L_{\text{\itshape Ter}}}
\newcommand{\LRel}{L_{\text{\itshape Rel}}}
\newcommand{\LFor}{L_{\text{\itshape For}}}
\newcommand{\NTer}{N_{\text{\itshape Ter}}}
\newcommand{\NRel}{N_{\text{\itshape Rel}}}
\newcommand{\NFor}{N_{\text{\itshape For}}}
\newcommand{\PTer}{P_{\text{\itshape Ter}}}
\newcommand{\PRel}{P_{\text{\itshape Rel}}}
\newcommand{\PFor}{P_{\text{\itshape For}}}

\newcommand{\plka}{\alka}
\newcommand{\plkz}{\alkz}
%\newcommand{\plka}{\plfoo{(}}
%\newcommand{\plkz}{\plfoo{)}}
\newcommand{\plcomma}{\hbox{\texttt{\color{alcolor},}}}
\newcommand{\pleq}{{\color{alcolor}\,\dot=\,}}

% MODIFIED (DJ)
% previously: \newcommand{\plfoo}[1]{\mathtt{\color{alcolor}#1}}
\newcommand{\plfoo}[1]{\texttt{\color{alcolor}#1}}

\newcommand{\plc}{\plfoo{c}}
\newcommand{\pld}{\plfoo{d}}
\newcommand{\plf}{\plfoo{f}}
\newcommand{\plg}{\plfoo{g}}
\newcommand{\plh}{\plfoo{h}}
\newcommand{\plx}{\plfoo{x}}
\newcommand{\ply}{\plfoo{y}}
\newcommand{\plz}{\plfoo{z}}
\newcommand{\plR}{\plfoo{R}}
\newcommand{\plS}{\plfoo{S}}

\newcommand{\bv}{\mathrm{bv}}
\newcommand{\fv}{\mathrm{fv}}

%\newcommand{\AxAL}{\hbox{\textit{Ax}}_{AL}}
%\newcommand{\AALV}{\hbox{\textit{Var}}_{AL}}

%\renewcommand{\#}[1]{\literal{#1}}
\newcommand{\A}{\mathcal{A}}
\newcommand{\Adr}{\text{Adr}}
\newcommand{\ar}{\mathrm{ar}}
\newcommand{\ascii}[1]{\literal{\char#1}}
%\newcommand{\assert}[1]{\text{/\!\!/\ } #1}
\newcommand{\assert}[1]{\colorbox{black!7!white}{\ensuremath{\{\;#1\;\}}}}
\newcommand{\Assert}[1]{$\langle$\textit{#1}$\rangle$}
\newcommand{\B}{\mathcal{B}}
\newcommand{\bfmod}{\mathbin{\kw{ mod }}}
\newcommand{\bb}{{\text{bb}}}
\def\bottom{\hbox{\small$\pmb{\bot}$}}
\newcommand{\card}[1]{|#1|}
%\newcommand{\cod}{\mathop{\text{cod}}}  % ist in thwmathabbrevs
\newcommand{\Conf}{\mathcal{C}}
\newcommand{\define}[1]{\emph{#1}}
%\renewcommand{\dh}{d.\,h.\@\xspace}
%\newcommand{\Dh}{D.\,h.\@\xspace}
%\newcommand{\engl}[1]{engl.\xspace\emph{#1}}
\newcommand{\eps}{\varepsilon}
%\newcommand{\evtl}{evtl.\@\xspace}
\newcommand{\fbin}{\text{bin}}
\newcommand{\finv}{\text{inv}}
\newcommand{\fnum}{\text{num}}
\newcommand{\fNum}{{\text{Num}}}
\newcommand{\frepr}{\text{repr}}
\newcommand{\fRepr}{\text{Repr}}
\newcommand{\fZkpl}{\text{Zkpl}}
\newcommand{\fLen}{\text{Len}}
\newcommand{\fsem}{\text{sem}}
\providecommand{\fspace}{\mathord{\text{space}}}
\providecommand{\fSpace}{\mathord{\text{Space}}}
\providecommand{\ftime}{\mathord{\text{time}}}
\providecommand{\fTime}{\mathord{\text{Time}}}
\newcommand{\fTrans}{\text{Trans}}
\newcommand{\fVal}{\text{Val}}

% MODIFIED (DJ)
\newcommand{\Val}{\text{Val}}

%\def\G{\mathbb{Z}}
\newcommand{\HT}[1]{\normalfont\textsc{HT-#1}}
\newcommand{\htr}[3]{\{#1\}\;#2\; \{#3\}}
\newcommand{\Id}{\text{I}}
%\newcommand{\ie}{i.\,e.\@\xspace}
\newcommand{\instr}[2]{\texttt{#1}\ \textit{#2}}
\newcommand{\Instr}[2]{\texttt{#1}\ \textrm{#2}}
\newcommand{\instrr}[3]{\texttt{#1}\ \textit{#2}\texttt{(#3)}}
\newcommand{\Instrr}[3]{\texttt{#1}\ \textrm{#2}\texttt{(#3)}}
\newcommand{\io}{\!\mid\!}
\usepackage{KITcolors}
\newcommand{\literal}[1]{\hbox{\textcolor{blue!95!white}{\textup{\texttt{\scalebox{1.11}{#1}}}}}}
%\newcommand{\literal}[1]{\hbox{\textcolor{KITblue!80!black}{\textup{\texttt{#1}}}}}
\def\kasten#1{\leavevmode\literal{\setlength{\fboxsep}{1pt}\fbox{\vrule  width 0pt height 1.5ex depth 0.5ex #1}}}
\newcommand{\kw}[1]{\ensuremath{\mathbf{#1}}}
\newcommand{\lang}[1]{\ensuremath{\langle#1\rangle}}
%\newcommand{\maw}{m.\,a.\,w.\@\xspace}
%\newcommand{\MaW}{M.\,a.\,w.\@\xspace}
\newcommand{\mdefine}[2][FOOBAR]{\define{#2}\def\foobar{FOOBAR}\def\optarg{#1}\ifx\foobar\optarg\def\optarg{#2}\fi\graffito{\optarg}}
\newcommand{\meins}{\rotatebox[origin=c]{180}{1}}
\newcommand{\Mem}{\text{Mem}}
\newcommand{\memread}{\text{memread}}
\newcommand{\memwrite}{\text{memwrite}}
\providecommand{\meta}[1]{\ensuremath{\langle}\textit{#1}\ensuremath{\rangle}}
%\newcommand{\N}{\mathbb{N}}
\newcommand{\NP}{\mathbf{NP}}
\newcommand{\Nadd}{N_{\text{add}}}
\newcommand{\Nmult}{N_{\text{mult}}}
% MODIFIED (DJ): added \!, mathcal{O}
\newcommand{\Oh}[1]{\mathcal{O}\!\left(#1\right)}
\newcommand{\Om}[1]{\Omega\!\left(#1\right)}
\newcommand{\personname}[1]{\textsc{#1}}
\newcommand{\regname}[1]{\texttt{#1}}
\newcommand{\mima}{\textsc{Mima}\xspace}
\newcommand{\mimax}{\textsc{Mima-X}\xspace}

\def\Pclass{\text{\bfseries P}}
\def\PSPACE{\text{\bfseries PSPACE}}

\newcommand{\SPush}{\text{push}}
\newcommand{\SPop}{\text{pop}}
\newcommand{\SPeek}{\text{peek}}
\newcommand{\STop}{\text{top}}
\newcommand{\STos}{\text{\itshape tos}}
\newcommand{\SBos}{\text{\itshape bos}}

%\newcommand{\R}{\mathbb{R}}
\newcommand{\Rnullplus}{\R_0^{+}}
\newcommand{\Rplus}{\R_{+}}
\newcommand{\resp}{resp.\@\xspace}
\newcommand{\Sem}{\text{Sem}}
\newcommand{\sgn}{\mathop{\text{sgn}}}
\newcommand{\sqbox}{\mathop{\raisebox{-6.2pt}{\hbox{\hbox to 0pt{$^{^{\sqcap}}$\hss}$^{^{\sqcup}}$}}}}
\newcommand{\sqleq}{\sqsubseteq}
\newcommand{\sqgeq}{\sqsupseteq}
% MODIFIED (DJ): added \!
\newcommand{\Th}[1]{\Theta\!\left(#1\right)}
%\newcommand{\usw}{usw.\@\xspace}
\newcommand{\V}[1]{\hbox{\textit{#1}}}
\newcommand{\x}{\times}
\newcommand{\ZK}{\mathbb{K}}
%\newcommand{\Z}{\mathbb{Z}}
\newcommand{\zB}{z.\,B.\@\xspace}
\newcommand{\ZB}{Z.\,B.\@\xspace}
% \newcommand{\bb}{{\text{bb}}}
% \def\##1{\hbox{\textcolor{darkblue}{\texttt{#1}}}}
% \def\A{\mathcal{A}}
% \newcommand{\0}{\#0}
% \newcommand{\1}{\#1}
% \newcommand{\Obj}{\text{Obj}}
% \newcommand{\start}{\mathop{\text{start}}}
% \newcommand{\compactlist}{\addtolength{\itemsep}{-\parskip}}
% \newcommand{\fval}{\text{val}}
% \newcommand{\lang}[1]{\ensuremath{\langle#1\rangle}}
% \newcommand{\io}{\!\mid\!}
% \def\sqbox{\mathop{\raisebox{-6.2pt}{\hbox{\hbox to 0pt{$^{^{\sqcap}}$\hss}$^{^{\sqcup}}$}}}}
% \def\sqleq{\sqsubseteq}
% \def\sqgeq{\sqsupseteq}
\def\Td{T_{\overline{d}}}
% \newcommand{\csym}[1]{\ensuremath{\#{c}_{\#{\hbox{\scriptsize #1}}}}}
% \newcommand{\F}{\ensuremath{\mathcal{F}}}
% \newcommand{\fsym}[2]{\ensuremath{\#{f}^{\#{\hbox{\scriptsize #1}}}_{\#{\hbox{\scriptsize #2}}}}}
% \newcommand{\rsym}[2]{\ensuremath{\#{R}^{\#{\hbox{\scriptsize #1}}}_{\#{\hbox{\scriptsize #2}}}}}
% \newcommand{\xsym}[1]{\ensuremath{\#{x}_{\#{\hbox{\scriptsize #1}}}}}
% \newcommand{\I}{\mathcal{I}}
% ********************************************************************

\usepackage{../TutTexbib/thwregex}

\newcommand{\nM}{\mathbb{M}}
\newcommand{\nR}{\mathbb{R}}
\newcommand{\nN}{\mathbb{N}}
\newcommand{\nZ}{\mathbb{Z}}
\newcommand{\nQ}{\mathbb{Q}}
\newcommand{\nB}{\mathbb{B}}
\newcommand{\nC}{\mathbb{C}}
\newcommand{\nK}{\mathbb{K}}
\newcommand{\nF}{\mathbb{F}}
\newcommand{\nG}{\mathbb{G}}
\newcommand{\nullel}{\mathcal{O}}
\newcommand{\einsel}{\mathds{1}}
\newcommand{\nP}{\mathbb{P}}
\newcommand{\Pot}{\mathcal{P}}
\renewcommand{\O}{\text{O}}

%\newcommand{\literal}[1]{\textcolor{blue}{\texttt{#1}}}
\newcommand{\realTilde}{\textasciitilde \ }
\newcommand{\setsize}[1]{\; \mid #1 \mid \; }

\definecolor{myRed}{RGB}{255,75,20}
\colorlet{myGreen}{KITpalegreen}

\newcounter{tfqtempcount}
\newcommand{\truefalseQ}[4]{
	\setcounter{tfqtempcount}{#1}
	\addtocounter{tfqtempcount}{1}
	\truefalseQuestion{#1}{\value{tfqtempcount}}{#2}{#3}{#4}
}
\newcommand{\truefalseQuestion}[5]{\item<#1-|handout:#1-> \color<#2-|handout:#2->{#3} #4 \qquad \visible<#2-|handout:#2->{#5}}

\newcommand{\boder}{\ensuremath{\textcolor{blue}{\vee}}\text{} }
\newcommand{\bund}{\ensuremath{\textcolor{blue}{\wedge}}\text{}}
\newcommand{\bimp}{\ensuremath{\textcolor{blue}{\to}}\text{}}
\newcommand{\bnot}{\ensuremath{\textcolor{blue}{\neg}}\text{}}
\newcommand{\bone}{\ensuremath{\textcolor{blue}{1}}\text{}}
\newcommand{\bzero}{\ensuremath{\textcolor{blue}{0}}\text{}}
\newcommand{\bleftBr}{\ensuremath{\textcolor{blue}{(}}\text{}}
\newcommand{\brightBr}{\ensuremath{\textcolor{blue}{)}}\text{}}

\newcommand{\summe}[2]{\sum\limits_{#1}^{#2}}
\newcommand{\limes}[1]{\lim\limits_{#1}}

\newcommand{\numpp}{\advance \value{weeknum} by -2 \theweeknum \advance \value{weeknum} by 2}
\newcommand{\nump}{\advance \value{weeknum} by -1 \theweeknum \advance \value{weeknum} by 1}

\newcommand{\mycomment}[1]{}

\newcounter{abc}
\newenvironment{alist}{
  \begin{list}{(\alph{abc})}{
      \usecounter{abc}\setlength{\leftmargin}{8mm}\setlength{\labelsep}{2mm}
    }
}{\end{list}}

\renewcommand{\mod}{\text{ \bf{mod} }}
\renewcommand{\div}{\text{ \bf{div} }}

%\newcommand{\sgn}{\text{sgn}}

% Toggle Handout mode by including this line before including style_tut:
% \newcommand{\handout}{}

%% Beamer-Klasse im korrekten Modus
\ifdefined \handout
\documentclass[handout]{beamer} % Handout mode
\else
\documentclass{beamer}
\fi


% Das ist der KIT-Stil
%\usepackage{../TutTexbib/beamerthemekit}
\usepackage[deutsch,titlepage0]{../TutTexbib/KIT/beamerthemeKITmod}
\TitleImage[width=\titleimagewd]{../figures/titlepage.jpg}
%\usetheme[deutsch,titlepage0]{KIT}

% Include PDFs
\usepackage{pdfpages}

% Libertine font (Original GBI font)
\usepackage{libertine}
%\renewcommand*\familydefault{\sfdefault}  %% Only if the base font of the document is to be sans serif

% Nicer math symbols
\usepackage{eulervm}
%\usepackage{mathpazo}

%% Deutsche Silbentrennung und Beschriftungen
\usepackage[ngerman]{babel}

%% UTF-8-Encoding
\usepackage[utf8]{inputenc}

\usepackage{csquotes}

%% Bibliotheken für viele mathematische Symbole
\usepackage{amsmath, amsfonts, amssymb}

%% Anzeigetiefe für Inhaltsverzeichnis: 1 Stufe
\setcounter{tocdepth}{1}

%% Schönere Schriften
\usepackage[TS1,T1]{fontenc}

%% Tabellen
\usepackage{array}
\usepackage{multicol}

%% Bibliothek für Graphiken
\usepackage{graphicx}

%% der wird sowieso in jeder Datei gesetzt
%%\graphicspath{{../figures/}}

%% Hyperlinks
\usepackage{hyperref}
% I don't know why, but this works and only includes sections and NOT subsections in the pdf-bookmarks.
\hypersetup{bookmarksdepth=subsection} 

%\usepackage{lmodern}
\usepackage{colortbl}
\usepackage[absolute,overlay]{textpos}
\usepackage{listings}
\usepackage{forloop}
%\usepackage{algorithmic} % PseudoCode package 

\usepackage{tikz}
\usetikzlibrary{matrix}
\usetikzlibrary{arrows.meta}
\usetikzlibrary{automata}
\usetikzlibrary{tikzmark}

% Needed for gbi-macros
\usepackage{xspace}

%%%%%%%%%%%% INHALT %%%%%%%%%%%%%%%%

%% Wochennummer
\newcounter{weeknum}

%% Titelinformationen
\title[GBI Tutorium, Woche \theweeknum]{Grundbegriffe der Informatik \\ Tutorium 36}
\subtitle{Termin \theweeknum \ | \mydate \\ Thassilo Helmold}
\author[Thassilo Helmold]{Thassilo Helmold}
\institute{KIT -- Karlsruher Institut für Technologie}
\date{\mydate}

%% Titel einfügen
\newcommand{\titleframe}{\frame{\titlepage}}

%% Alles starten mit \starttut{X}
\newcommand{\starttut}[1]{\setcounter{weeknum}{#1}\titleframe\frame{\frametitle{Inhalt}\tableofcontents} \AtBeginSection[]{%
\begin{frame}
	\tableofcontents[currentsection]
\end{frame}\addtocounter{framenumber}{-1}}}


\newcommand{\framePrevEpisode}{
	\begin{frame}
		\centering
		\textbf{In the previous episode of GBI...}
	\end{frame}
}

%% Legacy: Lastframe. Not for further usage!
\newcommand{\lastframe}[4]{\xkcdframe{#1}{#2}{#3}{#4}{}}
\newcommand{\xkcdframe}[5]{
	\frame[plain]{
		\vspace{-#2pt}
		\begin{figure}[H]
			\centering
			\includegraphics[scale=#1]{#3}
			\vspace{-7pt}
			\caption{ \texttt{\url{#4}} }
			\vspace{5pt}
			#5
		\end{figure} 
	}
}

%% Wörter
\newcommand{\code}[1]{$\mathbf{#1}$}

%% Sterne

\newcounter{starsc}
\newcommand{\stars}[1]{
	\hfill
	\begin{minipage}{100px}
		\forloop{starsc}{0}{\value{starsc} < #1}%
		{%
			\includegraphics[scale=0.05]{star-full.pdf} \hspace*{1px}
		}%
		\forloop{starsc}{\value{starsc}}{\value{starsc} < 5}%
		{%
			\includegraphics[scale=0.05]{star-empty.pdf} \hspace*{1px}
		}
		\vspace*{2px}
	\end{minipage}
}

% No stars for me...
\renewcommand{\stars}[1]{}

\newcommand{\slideThanks}{
	\begin{frame}
		\frametitle{Danksagung}
		\begin{block}{}
			Dieser Foliensatz basiert in Teilen auf Folien von:\\[1em]
			Philipp Basler \\
			Nils Braun \\
			Dominik Doerner \\
			Ou Yue \\
		\end{block}
	\end{frame}
}

%% Verbatim
\usepackage{moreverb}

\usepackage{tikz}
\usetikzlibrary{arrows}
   \usetikzlibrary{automata}
   \usetikzlibrary{matrix}
\usepackage{url}

\graphicspath{{../figures/}}


\begin{document}
\starttut{8}

\lastframe{0.50}{0}{xkcd/logic_principle_of_explosion.png}{https://www.xkcd.com/}

\framePrevEpisode

\begin{frame}{Rückblick: MIMA}
	\begin{itemize}[<+->]
		\item Ein idealisierter Prozessor
		\item Einfach zu verstehen, aufwändig zu programmieren
		\item Hardware-Details beachten: Keine negativen Konstanten mit LDC möglich!
		\item Programme sind oftmals mit \enquote{Bit-Magie} einfacher und kürzer (aber auch schwerer zu verstehen)
	\end{itemize}
\end{frame}

\begin{frame}{Wahr oder Falsch?}
	\begin{block}{}
	\begin{itemize}
		\FalseQuestionE{Im IR der MIMA wird die Adresse des aktuellen Befehls gespeichert}{Das geschieht im IAR. Im IR steht der Befehl selbst}
		\TrueQuestionE{Die Befehlsholphase ist jede Ausführungsrunde identisch}{}
		\FalseQuestionE{Der Akku führt bei der MIMA Berechnungen aus}{Das macht die ALU. Im Akku wird das letzte Ergebnis zwischengespeichert}
	\end{itemize}
	\end{block}
\end{frame}

\begin{frame}{Wahr oder Falsch?}
	\begin{block}{}
	\begin{itemize}
	\FalseQuestionE{LDC -5 lädt $-5$ in den Akku}{LDC funktioniert nicht mit negativen Konstanten!}
	\FalseQuestionE{Mit der MIMA können wir Zufallszahlen erzeugen}{Die MIMA arbeitet rein deterministisch und damit ohne jeden Zufall.}
	\end{itemize}
	\visible<5>{
	\begin{figure}[H]
		\centering
		\includegraphics[scale=0.6]{xkcd/random_number}
		\vspace{-7pt}
		\caption{ \texttt{\url{http://xkcd.com/}} }
	\end{figure} 	
	}
	\end{block}
\end{frame}

%% Zum Aufwärmen

\begin{frame}{Der Philosoph am Scheideweg}
	\only<1|handout:1>{
	Ein Weg führt in die Wüste, auf dem ein Wanderer marschiert. 
	Nach einem eintägigen Fußmarsch gelangt er an eine Weggabelung, an der eine kleine Hütte steht.\\
	\bigskip
	Neben der Eingangstüre ist ein Plakat mit folgendem Text befestigt:
	}
	\only<2|handout:1>{
	Einer der beiden weiterführenden Wege führt zu mehreren guten Oasen; nur auf diesem Weg kannst du dein Ziel lebend erreichen.\\
	\bigskip
	Der andere Weg führt in die unendliche Wüste; er hat noch jedem den Tod gebracht, der ihn wählte.\\
	\bigskip
	}
	\only<2-3|handout:2>{
	In dieser Hütte haust ein Philosoph, der die Ziele der beiden Wege sicher kennt; \textbf{streng abwechselnd sagt er an einem Tag die Wahrheit und am darauf folgenden Tag lügt er}.\\
	\bigskip
	Keiner außer ihm weiß aber, ob der heute die Wahrheit sagt oder aber heute lügt.\\
	\bigskip
	}
	\only<3|handout:2>{
	Du darfst dem Philosophen nur eine einzige Frage stellen, mit der du den Weg zu den Oasen erfahren willst.\\
	\bigskip
	Was fragst du ihn?
	}
\end{frame}

\begin{frame}{Der Philosoph am Scheideweg}
	Was fragst du ihn?\\[1em] \pause
	Was würdest du sagen, wenn ich dich morgen fragen würde, ob der linke Weg der richtige ist?\\
	\bigskip
	%Und wenn der Philosoph jeden Morgen würfelt, ob er die Wahrheit sagt oder lügt?\\ 
	%\only<2|handout:2>{
	%Was würdest du sagen, wenn ich dich \textbf{heute} fragen würde, ob der linke Weg der richtige ist?\\ }
\end{frame}

\begin{frame}{Rückblick: Kontextfreie Grammatiken}
	\begin{itemize}[<+->]
		\item Ein Vier-Tupel: $G = (N, T, S, P)$
		\item Produktionen definieren Ersetzungen eines Nichtterminals mit Wörtern über $N \cup T$
		\item Wir wenden Produktionen in Ableitungsschritten an: $v \Rightarrow w, v \Rightarrow^* w$
		\item $L(G)$ sind alle aus $S$ ableitbaren Wörter über $T$ (die also nur aus Terminalsymbolen bestehen)
	\end{itemize}
\end{frame}

\begin{frame}{Wahr oder Falsch?}
	Sei $G=(\{X,Y\},\{a,b\},X,P)$ eine kontextfreie Grammatik.
	\begin{itemize}
		\FalseQuestion{Die Produktion $XY \to a$ ist gültig}
		\FalseQuestion{Die Produktion $a \to XY$ ist gültig}
		\TrueQuestionE{Die Produktion $X \to XaX$ ist gültig}{}
		\TrueQuestionE{Wenn $X \to w$ eine gültige Produktion ist, dann gilt $X \Rightarrow w$}{}
		\FalseQuestionE{Wenn $X \Rightarrow w$ gilt, dann ist $X \to w$ eine Produktion in P }{Sei $P=\{X \to XX \mid a\}$. Dann gilt $XX \Rightarrow Xa$, aber $XX \to Xa \notin P$}	\end{itemize}
\end{frame}

\mycomment{ % Kam letztes Jahr nicht gut an, ist wohl etwas "albern"
\begin{frame}{Zum Aufwärmen: Weihnachtsgebäck}
	Weihnachten steht vor der Tür, aber Geschenke sind von Jahr zu Jahr teurer geworden! Um den Staatshaushalt des Nordpols zu schonen, hat sich der Weihnachtsmann etwas einfallen lassen: Dieses Jahr bringt er nur den Haushalten Geschenke, bei denen das für ihn bereitgestellte Gebäck seinen speziellen Qualitätsanforderungen genügt.\\[1em]
	
	Dazu hat der Weihnachtsmann drei Gebäckarten (Lebkuchen $L$, Zimtstern $Z$, Dominostein $D$) normiert und eine kontextfreie Grammatik aufgestellt, die gültige Kombinationen beschreibt.
\end{frame}

\begin{frame}{Zum Aufwärmen: Weihnachtsgebäck}
	\begin{align*}
		G &= (\{S, A, B, X, Y\}, \{L, Z, D\}, S, \{ \\
		S &\to A \mid B\\
		A &\to LLLLLLLLLLX, \qquad X \to LX \mid e\\
		B &\to ZYZ \mid DYD, \qquad Y \to Z \mid D \mid L\\
		\})
	\end{align*}
	\begin{enumerate}
		\item Welche Kombinationen sind gültig?\\[0.5em]
		\item Der Weihnachtsmann muss auf seine Figur achten. Deswegen möchte er für nächstes Jahr eine neue, kalorienarme Kombination von Gebäck:\\
		3 Folgen mit 42, 17 und 41 x ein Zimtstern oder ein Dominostein. Die einzelnen Folgen sollen durch Lebkuchen getrennt werden.\\
		Stelle eine Kontextfreie Grammatik dafür auf.
	\end{enumerate}
\end{frame}

\begin{frame}{Zum Aufwärmen: Weihnachtsgebäck}
	\begin{align*}
		G &= (\{S, A, B, X, Y\}, \{L, Z, D\}, S, \{ \\
		S &\to A \mid B\\
		A &\to LLLLLLLLLLX, \qquad X \to LX \mid e\\
		B &\to ZYZ \mid DYD, \qquad Y \to Z \mid D \mid L\\
		\})
	\end{align*}

	Welche Kombinationen sind gültig?\\[0.5em] \pause
	$\{ L^n \mid n \geq 10\} \cup \{w \mid \text{w ist Palindrom und } |w|_L \leq 1\}$
\end{frame}

\begin{frame}{Zum Aufwärmen: Weihnachtsgebäck}
	Der Weihnachtsmann muss auf seine Figur achten. Deswegen möchte er für nächstes Jahr eine neue, kalorienarme Kombination von Gebäck:\\
	3 Folgen mit 42, 17 und 41 x ein Zimtstern oder ein Dominostein. Die einzelnen Folgen sollen durch Lebkuchen getrennt werden.\\
	Stelle eine Kontextfreie Grammatik dafür auf.
	
	\pause
	\begin{align*}
		G &= (\{S\} \cup \{F_i \mid 0 \leq i \leq 42 \}, \{L, Z, D\}, S, P)\\
		P &= \{S \to F_{42}LF_{17}LF_{41} \}\\
		  &\cup \{F_i \to ZF_{i-1}, F_i \to DF_{i-1} \mid 1 \leq i \leq 42 \} \\
		  &\cup \{F_0 \to \eps \}
	\end{align*}
\end{frame}
}

\newcommand{\plB}{\plfoo{B}}
\newcommand{\plE}{\plfoo{E}}

\begin{frame}{}
	\begin{block} {Aussagenlogik}
	\begin{itemize}
		\item \enquote{Es regnet und alle Vögel sind grau.}
		\item atomar: \enquote{Es regnet.}, \enquote{ Alle Vögel sind grau.}
		\item Diese beiden Aussagen lassen sich ihrerseits nicht in weitere Teilaussagen zerlegen!
	\end{itemize}
	\end{block}

	\pause
	\begin{block} {Prädikatenlogik}	
		\begin{itemize}
		\item In der Prädikatenlogik werden atomare Aussagen hinsichtlich ihrer inneren Struktur untersucht.
		\item \enquote{Alle Vögel sind grau}
		\item lässt sich in : \enquote{Alle Vögel}, \enquote{sind grau} zerlegen.
	\end{itemize}
	
	\end{block}
\end{frame}

\section{Prädikatenlogik: Syntax}

\begin{frame}{Aufgabe 1 (15/16, Blatt 7)}
	  Es seien $\CPL = \{ \}$, $\VPL = \{ \plx, \ply, \plz \}$, $\FPL = \{ \}$ und $\RPL = \{ \plE, \pleq \}$ mit $\ar(\plE) = 2$, und es sei $F$ die prädikatenlogische Formel
	\begin{equation*}
	\alnot \plexist \plx
	{\plka
		\plE{\plka \plx \plcomma \ply \plkz}
		\alor
		\alnot \plall \plz \plall \plx \plall \ply
		{\plka
			\plE{\plka \plx \plcomma \plz \plkz} \aland \plE{\plka \ply \plcomma \plz \plkz} \alimpl \plx \pleq \ply
			\plkz}
		\plkz}
	\end{equation*}
	
	\begin{block}{Aufgabe 1.0}
	Ist diese prädikatenlogische Formel syntaktisch korrekt?\\ \pause
	Ja
	\end{block}
\end{frame}

\begin{frame}{Syntax}
	\begin{block}{Aufbau von prädikatenlogischen Formeln}
	\begin{itemize}[<+->]
		\item \textbf{Terme}: Liefern \enquote{Werte}; Aus Konstanten, Variablen und Funktionssymbolen zusammengesetzt.
		\item \textbf{Atomare Formeln}: Liefern \enquote{Wahrheitswerte}; Aus Termen und Relationssymbolen zusammengesetzt.
		\item \textbf{Prädikatenlogische Formeln}: Verknüpfen und Quantifizieren atomare Formeln; Aus atomaren Formeln und aussagenlogischen Konnektiven sowie Quantoren zusammengesetzt. 
	\end{itemize}
	\end{block}
\end{frame}


\begin{frame}{Terme - Alphabet}
\begin{figure}
	\centering
	\includegraphics[scale=0.2]{pl/TermeAlphabet.png} \hspace{2em} 
\end{figure} 
\end{frame}

\begin{frame}{Terme - Grammatik}
\begin{figure}[h!]
	\centering
	\includegraphics[scale=0.2]{pl/TermeGrammatik.png} \hspace{2em} 
\end{figure} 
\end{frame}

\begin{frame}{Atomare Formeln}
\begin{figure}[h!]
	\centering
	\includegraphics[scale=0.2]{pl/AFormeln.png} \hspace{2em} 
\end{figure} 
\end{frame}

\begin{frame}{Atomare Formeln - Beispiel}
\begin{figure}[h!]
	\centering
	\includegraphics[scale=0.2]{pl/AFormelnBsp.png} \hspace{2em} 
\end{figure} 
\end{frame}

\begin{frame}{Prädikatenlogische Formeln}
\begin{figure}[h!]
	\centering
	\includegraphics[scale=0.2]{pl/PFormeln.png} \hspace{2em} 
\end{figure} 
\end{frame}

\begin{frame}{Quiz}
\begin{figure}[h!]
	\centering
	\only<1|handout:1>{\includegraphics[scale=0.2]{pl/q1.png}}
	\only<2|handout:2>{\includegraphics[scale=0.2]{pl/q2.png}} 
	\only<3|handout:3>{\includegraphics[scale=0.2]{pl/q3.png}}
	\only<4|handout:4>{\includegraphics[scale=0.2]{pl/q4.png}}
	\only<5|handout:5>{\includegraphics[scale=0.2]{pl/q5.png}}
	\only<6|handout:6>{\includegraphics[scale=0.2]{pl/q6.png}}
	\only<7|handout:7>{\includegraphics[scale=0.2]{pl/q7.png}}
	\hspace{2em}
\end{figure} 
\end{frame}

\begin{frame}{Freie und gebundene Variablenvorkommen}
	\begin{equation*}
	F = \alnot \plexist \plx
	{\plka
		\plE{\plka \plx \plcomma \ply \plkz}
		\alor
		\alnot \plall \plz \plall \plx \plall \ply
		{\plka
			\plE{\plka \plx \plcomma \plz \plkz} \aland \plE{\plka \ply \plcomma \plz \plkz} \alimpl \plx \pleq \ply
			\plkz}
		\plkz}
	\end{equation*}
	
	\begin{block}{Aufgabe 1.1}
		Welche Variablenvorkommen sind frei ($\fv$) und welche gebunden ($\bv$)?\\
		Ist die Formel geschlossen?
	\end{block}

	\pause
	\begin{block}{Lösung}
		Nur die Variable $\fv(F) = \{\ply\}$ kommt frei in $F$ vor.\\
		Genau die Variablen $\bv(F) = \{\plx, \ply, \plz\}$ kommen gebunden in $F$ vor.\\
		Da $\fv(F) \neq \emptyset$ ist $F$ nicht geschlossen.
	\end{block}
	
\end{frame}


\begin{frame}{Substitutionen}
	\begin{figure}[h!]
		\centering
		\includegraphics[scale=0.6]{pl/subs.png} \hspace{2em} 
	\end{figure} 
\end{frame}

\begin{frame}{Substitutionen}
	Ersetzt werden nur \textbf{freie Variablenvorkommen}!\\
	Gebundene Vorkommen, also Variablen im Wirkungsbereich eines Quantors, werden \textbf{nicht} ersetzt.
	
	\pause
	\begin{figure}[h!]
		\centering
		\includegraphics[scale=0.5]{pl/subs_bsp.png} \hspace{2em} 
	\end{figure} 
\end{frame}

\begin{frame}{Substitutionen: Kollisionsfreiheit}
	Bei einer \textbf{kollisionsfreien} Substitution werden keine Variablen \enquote{aus Versehen} gebunden.\\[1em]
	Ersetzen wir eine freie Variable $x$ durch einen Term, in dem die Variable $y$ frei vorkommt, so darf sich $x$ nicht im Wirkungsbereich eines Quantors über $y$ befinden.
	
	\pause
	\begin{Beispiel}
		$L = \plall \plx (\plx \aland \ply)$\\
		Kollisionsfrei: $\sigma_{\{\ply/\plz\}}$\\
		Nicht kollisionsfrei: $\sigma_{\{\ply/\plx\}}$
	\end{Beispiel}
\end{frame}

\begin{frame}{Substitutionen}
	\begin{equation*}
	F = \alnot \plexist \plx
	{\plka
		\plE{\plka \plx \plcomma \ply \plkz}
		\alor
		\alnot \plall \plz \plall \plx \plall \ply
		{\plka
			\plE{\plka \plx \plcomma \plz \plkz} \aland \plE{\plka \ply \plcomma \plz \plkz} \alimpl \plx \pleq \ply
			\plkz}
		\plkz}
	\end{equation*}
	
	\begin{block}{Aufgabe 1.2}
		Geben sie eine Substitution $\sigma$ an, die \emph{nicht} kollisionsfrei für $F$ ist.\\[1em] \pause
		
		Die Substitution $\sigma_{\{(\ply/\plx)\}}$ leistet das Gewünschte.
	\end{block}
	
\end{frame}


\section{Prädikatenlogik: Semantik}

\begin{frame}{Interpretation}
	\begin{figure}[h!]
		\centering
		\includegraphics[scale=0.6]{pl/int.png} \hspace{2em} 
	\end{figure} 
\end{frame}

\begin{frame}{Aufgabe 1 (15/16, Blatt 7)}
	  \begin{equation*}
	\alnot \plexist \plx
	{\plka
		\plE{\plka \plx \plcomma \ply \plkz}
		\alor
		\alnot \plall \plz \plall \plx \plall \ply
		{\plka
			\plE{\plka \plx \plcomma \plz \plkz} \aland \plE{\plka \ply \plcomma \plz \plkz} \alimpl \plx \pleq \ply
			\plkz}
		\plkz}
	\end{equation*}\\[1em]
	
	\begin{block}{Aufgabe 1.3}
		Geben Sie eine Interpretation $(D_1, I_1)$ und eine Variablenbelegung $\beta_1$ so an, dass $val_{D_1, I_1, \beta_1}(F) = \textbf{W}$ gilt.\\[0.5em]
		
		\visible<2-|handout:2>{Die Interpretation $(D_1, I_1) = (\{ 0, 1 \}, {<})$ und die Variablenbelegung $\beta_1 \colon \VPL \to D$, $v \mapsto 0$, leisten das Gewünschte.}
	\end{block}

	\begin{block}{Aufgabe 1.4}
		Geben Sie eine Interpretation $(D_2, I_2)$ und eine Variablenbelegung $\beta_2$ so an, dass $val_{D_2, I_2, \beta_2}(F) = \textbf{F}$ gilt.\\[0.5em]
		
		\visible<3-|handout:2>{Die Interpretation $(D_2, I_2) = (\{ 0, 1 \}, {<})$ und die Variablenbelegung $\beta_2 \colon \VPL \to D$, $v \mapsto 1$, leisten das Gewünschte.}
	\end{block}
\end{frame}



\section{Prädikatenlogik: Aufgaben}

\begin{frame}{Prädikatenlogische Formeln aufstellen}
	% TODO Übung einbinden
	Vgl. Übung WS 15/16
\end{frame}

\begin{frame}{Aufgabe 2 (15/16, Blatt 7)}
	\begin{block}{Aufgabe}
		Formulieren Sie die folgenden Aussagen als Formeln in Prädikatenlogik:
		\begin{enumerate}
			\item Nicht alle Vögel können fliegen.
			\item Wenn es irgendjemand kann, dann kann es Donald Ervin Knuth.
			\item John liebt jeden, der sich nicht selbst liebt.
		\end{enumerate}
	\end{block}
	
	\only<2-|handout:2>{
	\begin{block}{Lösung}
		\begin{enumerate}
				\item
				\begin{equation*}
				\plexist \plx {\plka \plfoo{Vogel}{\plka \plx \plkz} \aland \alnot \plfoo{flugfaehig}{\plka \plx \plkz} \plkz}
				\end{equation*}
			\only<3-|handout:2>{
				\item
				\begin{equation*}
				\plexist \plx {\plka \plfoo{kann\_es}{\plka \plx \plkz} \plkz}
				\alimpl
				\plfoo{kann\_es}{\plka \plfoo{knuth} \plkz}
				\end{equation*}
			}
			\only<4-|handout:2>{
				\item
				\begin{equation*}
				\plall \plx {\plka \alnot \plfoo{liebt}{\plka \plx \plcomma \plx \plkz} \alimpl \plfoo{liebt}{\plka \plfoo{John} \plcomma \plx \plkz} \plkz}
				\end{equation*}
			}
		\end{enumerate}
	\end{block}
	}
\end{frame}

\begin{frame}{Weitere Aufgaben}
	Siehe Übung WS 15/16
\end{frame}

% Evtl. Übung: Nicht-Palindrome?

\begin{frame}	
	\begin{block}{Was ihr nun wissen solltet}
		\begin{itemize}
			\item Wie Prädikatenlogische Formeln aufgebaut sind...
			\item ... und was sie bedeuten.
		\end{itemize}
	\end{block}
	
	\begin{block}{Was nächstes Mal kommt}
		\begin{itemize}
			\item Alles Korrekt? - Beweise mit dem Hoare-Kalkül
		\end{itemize}
	\end{block}
\end{frame}

\begin{frame}[plain]
	\begin{center}
		\large
		Vergesst nicht: Nächste Woche findet noch ein Tutorium statt!
	\end{center}
	\bigskip
	Für alle die nicht kommen: \\
	Frohe Weihnachten und einen guten Start in das neue Jahr!
\end{frame}

\lastframe{0.50}{0}{xkcd/christmastree.png}{https://www.xkcd.com/835}
\slideThanks

\end{document}